% GENERATED FILE. Edit canonical lessons, then run npm run generate:books.
% canonical-source-hash: fnv1a64:44e7fe1a5edf6f87
% canonical-lessons: ZH-C287-dang, ZH-W287-dang, ZH-C287-ambition, ZH-C287-mengxiang, ZH-C287-mubiao, ZH-C287-zui-most, ZH-C287-practice

\chapter{What I Want to Be}
\label{ch:zh-what-i-want-to-be}
\hlchaptermodality{287}

\begin{chapteropening}
\textbf{By the end of this chapter:} \emph{I can say what I want to do or be later.}

Saying what you want to do or be with \zh{以后我想} and \zh{当}, a dream and a goal, and \zh{最} before \zh{喜欢} or an adjective, built on one new written character: \zh{当}.
\end{chapteropening}

\section[\texorpdfstring{\zh{当} (dāng)}{当 (dāng)}]{\zh{当} (dāng) — to work as}

\label{lesson:ZH-C287-dang}



\begin{quote}
A new chapter, and it opens on the job you might do one day.

\begin{itemize}
\raggedright
  \item \emph{Recall:} say \zh{我感觉很好} — \emph{wǒ gǎnjué hěn hǎo} — \textbf{R1}, the chapter you have finished
  \item \emph{Recall:} say \zh{医生} — \emph{yīshēng}, a doctor — \textbf{R3}, durable
  \item \emph{Recall:} say \emph{lǎoshī}, a teacher — \textbf{R3}
  \item \emph{Recall:} say \emph{dāngshí}, at that time — \textbf{R3}
\end{itemize}
\end{quote}

\subsection*{You'll want to know: \zh{当}}
\textbf{\zh{当}} — \emph{dāng} — \textquotedblleft{}to work as; to be\textquotedblright{}, in a job or a role.

\begin{quote}
\textbf{\zh{我想当医生}} \emph{wǒ xiǎng dāng yīshēng} — I want to be a doctor.
\end{quote}

\begin{quote}
\emph{tā dāng le liǎng nián lǎoshī} — He worked as a teacher for two years.
\end{quote}

The next lesson puts it in your hand. It is the dāng of \emph{dāngshí}.

\begin{cousinweb}
\textbf{\zh{当}} (\emph{dāng}) is inherited from Middle Chinese \emph{tang} (Guangyun, in Baxter's transcription). The dictionary lists the character under two readings, \emph{tang} and \emph{tangH}, and glosses \emph{tang} 'to match', 'to face' and 'to be in charge'; 'to work as' is not given there, so how the word came by it is not settled.
\end{cousinweb}

\begin{grammarlens}[title={What you've built}]
\zh{当} takes a job or a role after it: \zh{当医生}, \emph{dāng lǎoshī}.

With a length of time, it follows the pattern you have: \emph{dāng le liǎng nián lǎoshī}.
\end{grammarlens}

\subsection*{Your turn}
\begin{itemize}
\raggedright
  \item \emph{Say it:} \emph{dāng}
  \item \emph{Say it:} \emph{dāng}, once more
  \item \emph{Say it:} \zh{我想当医生} — \emph{wǒ xiǎng dāng yīshēng} — I want to be a doctor.
  \item \emph{Say it:} \emph{tā dāng le liǎng nián lǎoshī} — He worked as a teacher for two years.
  \item \emph{Recall:} say \emph{dāngshí}, then \zh{当}: the same syllable
\end{itemize}

\begin{tcolorbox}[breakable,colback=teal!4,colframe=teal!35!black,title={Before you move on}]
What does \zh{当} mean? (\textquotedblleft{}To work as; to be\textquotedblright{}.) What follows it? (\textbf{A job or a role}.)
\end{tcolorbox}
\section[\texorpdfstring{\zh{当} (dāng)}{当 (dāng)}]{\zh{当} — the top of \zh{常}, then the foot of \zh{雪}}

\label{lesson:ZH-W287-dang}



\begin{quote}
The word first, then the character inside it.

\begin{itemize}
\raggedright
  \item \emph{Recall:} say \emph{dāng} — \textbf{R1}, one lesson back, before it decays
\end{itemize}

\begin{quote}
\zh{常} \zh{雪}
\end{quote}

\emph{Read it:} each one aloud, with nothing said first — \emph{cháng}, \emph{xuě}
\end{quote}

\subsection*{Script — the top of \zh{常}, then the foot of \zh{雪}}
\hlblockfigure{\detokenize{figures/ZH-W287-dang-filmstrip.pdf}}{How \zh{当} is written, stroke by stroke}

\begin{quote}
\zh{当}
\end{quote}

\textbf{Two pieces, and you write both already.}

\textbf{Three short strokes on top}, as on top of \zh{常}: the short vertical in the middle first, then the dot to its left, then the short left-falling stroke to its right.

\textbf{Three horizontals joined on the right}, as at the foot of \zh{雪}: a horizontal that turns down at its right end, a horizontal across the middle, and a horizontal along the foot.

3 + 3 = six.

\subsection*{Your turn}
\begin{itemize}
\raggedright
  \item \emph{Write it:} \zh{当}, all six, with the model in front of you
  \item \emph{Cover:} write it again from memory, and count the pieces, not the strokes
\end{itemize}

\begin{tcolorbox}[breakable,colback=teal!4,colframe=teal!35!black,title={Before you move on}]
How many strokes in \zh{当}? (\textbf{Six}.) Which stroke comes first? (\textbf{The short vertical in the middle}.) Where have you written its lower piece? (\textbf{At the foot of \zh{雪}}.)

Source: \href{https://raw.githubusercontent.com/chanind/hanzi-writer-data/68d10a4b21150cae5e1ebbd223eed289cf32d90c/data/\%E5\%BD\%93.json}{Hanzi Writer Data \zh{当}, pinned stroke paths and medians}.
\end{tcolorbox}
\section[\texorpdfstring{\zh{以后我想}… (yǐhòu wǒ xiǎng …)}{以后我想… (yǐhòu wǒ xiǎng …)}]{\zh{以后我想}… — later, I want to}

\label{lesson:ZH-C287-ambition}



\begin{quote}
Words you have, put together to look ahead.

\begin{itemize}
\raggedright
  \item \emph{Read it:} \zh{我想当医生}, aloud — \emph{wǒ xiǎng dāng yīshēng}
  \item \emph{Recall:} say \emph{jiānglái}, in the future — \textbf{R3}, durable
  \item \emph{Recall:} say \zh{以后} — \emph{yǐhòu}, later — \textbf{R3}
\end{itemize}
\end{quote}

\begin{grammarlens}[title={Grammar Lens: \zh{以后} and \zh{想}, looking ahead}]
\begin{quote}
\textbf{A time word for later, then \zh{想} or \zh{要}, says what you want to do or be.}
\end{quote}

\begin{quote}
\textbf{\zh{以后我想去中国工作}} \emph{yǐhòu wǒ xiǎng qù Zhōngguó gōngzuò} — Later on I want to go and work in China. \textbf{\zh{以后你想当什么}} \emph{yǐhòu nǐ xiǎng dāng shénme} — What do you want to be later on?
\end{quote}

\zh{以后} here is later, some time to come, at the front of the sentence. \emph{jiānglái} reaches further: some day, in the future.

\zh{想} is a wish; \zh{要} is a plan you mean to keep, as in \zh{下个月我要去中国}.
\end{grammarlens}

\subsection*{Your turn}
Say these aloud:

\begin{itemize}
\raggedright
  \item \zh{以后我想去中国工作} — \emph{yǐhòu wǒ xiǎng qù Zhōngguó gōngzuò} — Later on I want to go and work in China.
  \item \zh{以后你想当什么} — \emph{yǐhòu nǐ xiǎng dāng shénme} — What do you want to be later on?
  \item \emph{wǒ xiǎng dāng lǎoshī} — I want to be a teacher.
  \item what you want to do later on, with \zh{以后} and \zh{想}
\end{itemize}

\begin{tcolorbox}[breakable,colback=teal!4,colframe=teal!35!black,title={Before you move on}]
Where does \zh{以后} go? (\textbf{At the front}.) Which says a wish, \zh{想} or \zh{要}? (\textbf{\zh{想}}.)
\end{tcolorbox}
\section[mèngxiǎng]{mèngxiǎng — a dream for your life}

\label{lesson:ZH-C287-mengxiang}



\begin{quote}
Before the new one: last lesson's plan, and a word for what happens in sleep.

\begin{itemize}
\raggedright
  \item \emph{Read it:} \zh{以后我想去中国工作}, aloud — \emph{yǐhòu wǒ xiǎng qù Zhōngguó gōngzuò}
  \item \emph{Recall:} say \emph{zuòmèng}, to dream — \textbf{R3}, durable
\end{itemize}
\end{quote}

\subsection*{You'll want to know: mèngxiǎng}
\emph{mèngxiǎng} — \textquotedblleft{}a dream\textquotedblright{}, the one you have for your life.

\begin{quote}
\emph{wǒ de mèngxiǎng shì dāng yīshēng} — My dream is to be a doctor.
\end{quote}

Its first character comes later. The second is the \zh{想} you write.

\begin{cousinweb}
\emph{mèngxiǎng} is built from mèng 'a dream' and \zh{想} 'to think', both inherited from Middle Chinese. The Guangyun lists the first under two readings, \emph{mjuwng} and \emph{mjuwngH}, and under the second describes the spirit wandering in sleep; \zh{想} is \emph{sjangX}, glossed 'to think' (Guangyun, in Baxter's transcription).
\end{cousinweb}

\begin{grammarlens}[title={What you've built}]
\emph{zuòmèng} is dreaming in your sleep; \emph{mèngxiǎng} is the dream you have for your life.

It is a noun, and \zh{是} links it to the dream itself: \emph{wǒ de mèngxiǎng shì …}.
\end{grammarlens}

\subsection*{Your turn}
\begin{itemize}
\raggedright
  \item \emph{Say it:} \emph{mèngxiǎng}
  \item \emph{Say it:} \emph{mèngxiǎng}, once more
  \item \emph{Say it:} \emph{wǒ de mèngxiǎng shì dāng yīshēng} — My dream is to be a doctor.
  \item \emph{Recall:} say \emph{zuòmèng}, then \emph{mèngxiǎng}: asleep, then awake
\end{itemize}

\begin{tcolorbox}[breakable,colback=teal!4,colframe=teal!35!black,title={Before you move on}]
What does \emph{mèngxiǎng} mean? (\textquotedblleft{}A dream for your life\textquotedblright{}.) And \emph{zuòmèng}? (\textbf{To dream in your sleep}.)
\end{tcolorbox}
\section[mùbiāo]{mùbiāo — a goal}

\label{lesson:ZH-C287-mubiao}



\begin{quote}
Before the new one: last lesson's word, and one that shares this word's first sound.

\begin{itemize}
\raggedright
  \item \emph{Recall:} say \emph{mèngxiǎng}, a dream for your life — \textbf{R1}, one lesson back
  \item \emph{Recall:} say \emph{jiémù}, a programme — \textbf{R3}, durable
\end{itemize}
\end{quote}

\subsection*{You'll want to know: mùbiāo}
\emph{mùbiāo} — \textquotedblleft{}a goal; a target\textquotedblright{}.

\begin{quote}
\emph{wǒ de mùbiāo shì shuō hǎo Hànyǔ} — My goal is to speak Chinese well.
\end{quote}

Its characters come later; for now it is a word for the ear.

\begin{cousinweb}
\emph{mùbiāo} is built from mù 'the eye' and biāo 'a mark', both inherited from Middle Chinese, \emph{mjuwk} and \emph{pjiew} (Guangyun, in Baxter's transcription), which glosses the first 'the eye' and the second 'to raise' and 'the tip of a tree'. 'A goal' for the pair is its own sense.
\end{cousinweb}

\begin{grammarlens}[title={What you've built}]
\emph{mèngxiǎng} is what you hope for; \emph{mùbiāo} is what you aim at, and can reach step by step.
\end{grammarlens}

\subsection*{Your turn}
\begin{itemize}
\raggedright
  \item \emph{Say it:} \emph{mùbiāo}
  \item \emph{Say it:} \emph{mùbiāo}, once more
  \item \emph{Say it:} \emph{wǒ de mùbiāo shì shuō hǎo Hànyǔ} — My goal is to speak Chinese well.
  \item \emph{Recall:} say \emph{mèngxiǎng}, then \emph{mùbiāo}: the hope, then the aim
\end{itemize}

\begin{tcolorbox}[breakable,colback=teal!4,colframe=teal!35!black,title={Before you move on}]
What does \emph{mùbiāo} mean? (\textquotedblleft{}A goal\textquotedblright{}.) How is it different from \emph{mèngxiǎng}? (\textbf{A goal is something you aim at and can reach}.)
\end{tcolorbox}
\section[\texorpdfstring{\zh{我最喜欢}… (wǒ zuì xǐhuan …)}{我最喜欢… (wǒ zuì xǐhuan …)}]{\zh{我最喜欢}… — the one I like best}

\label{lesson:ZH-C287-zui-most}



\begin{quote}
A word you have said for a long time, printed for the first time.

\begin{itemize}
\raggedright
  \item \emph{Recall:} say \emph{mùbiāo}, a goal — \textbf{R1}, one lesson back
  \item \emph{Recall:} say \emph{zuì}, most — \textbf{R3}, durable
  \item \emph{Recall:} say \zh{他比我大} — \emph{tā bǐ wǒ dà} — \textbf{R3}
\end{itemize}
\end{quote}

\begin{grammarlens}[title={Grammar Lens: \zh{最}, the most of all}]
\begin{quote}
\textbf{\zh{最} before \zh{喜欢} or an adjective picks out the top of the group.}
\end{quote}

\begin{quote}
\textbf{\zh{我最喜欢中国菜}} \emph{wǒ zuì xǐhuan Zhōngguó cài} — I like Chinese food best. \textbf{\zh{这家书店最大}} \emph{zhè jiā shūdiàn zuì dà} — This bookshop is the biggest.
\end{quote}

\emph{zuì} you have said since the chapter on order; now it prints. \zh{比} sets two things side by side, \zh{我比他大}; \zh{最} needs no second thing: it is the most of all.

The story on the mountain ended on it: \zh{那天是我最开心的一天}.
\end{grammarlens}

\subsection*{Your turn}
Say these aloud:

\begin{itemize}
\raggedright
  \item \zh{我最喜欢中国菜} — \emph{wǒ zuì xǐhuan Zhōngguó cài} — I like Chinese food best.
  \item \zh{这家书店最大} — \emph{zhè jiā shūdiàn zuì dà} — This bookshop is the biggest.
  \item \zh{那天是我最开心的一天} — \emph{nà tiān shì wǒ zuì kāixīn de yì tiān} — That was my happiest day.
  \item the food you like best, with \zh{我最喜欢}
\end{itemize}

\begin{tcolorbox}[breakable,colback=teal!4,colframe=teal!35!black,title={Before you move on}]
Where does \zh{最} go? (\textbf{Before \zh{喜欢} or an adjective}.) Does it need a second thing, as \zh{比} does? (\textbf{No}.)
\end{tcolorbox}
\section[Practice]{\zh{以后我想当医生}}

\label{lesson:ZH-C287-practice}



\begin{quote}
The chapter closes here, with a question every child is asked.

\begin{itemize}
\raggedright
  \item \emph{Read it:} \zh{我最喜欢中国菜}, aloud — \emph{wǒ zuì xǐhuan Zhōngguó cài}
  \item \emph{Recall:} say \emph{mùbiāo}, a goal — \textbf{R2}
  \item \emph{Recall:} say \emph{mèngxiǎng}, a dream — \textbf{R2}
\end{itemize}
\end{quote}

\subsection*{The exchange}
\begin{quote}
\textbf{A:} \zh{以后你想当什么}
\end{quote}

\begin{quote}
\emph{yǐhòu nǐ xiǎng dāng shénme} — What do you want to be later on?
\end{quote}

\begin{quote}
\textbf{B:} \zh{我想当医生}
\end{quote}

\begin{quote}
\emph{wǒ xiǎng dāng yīshēng} — I want to be a doctor.
\end{quote}

\begin{quote}
\textbf{A:} \zh{为什么}
\end{quote}

\begin{quote}
\emph{wèishénme} — Why?
\end{quote}

\begin{quote}
\textbf{B:} \zh{因为我妈妈是医生}
\end{quote}

\begin{quote}
\emph{yīnwèi wǒ māma shì yīshēng} — Because my mother is a doctor.
\end{quote}

\begin{quote}
\textbf{A:} \emph{nǐ de mèngxiǎng shì shénme} — What's your dream?
\end{quote}

\begin{quote}
\textbf{B:} \emph{wǒ de mèngxiǎng shì qù Zhōngguó gōngzuò} — My dream is to work in China.
\end{quote}

\begin{quote}
\textbf{A:} \zh{你去过中国吗}
\end{quote}

\begin{quote}
\emph{nǐ qù guo Zhōngguó ma} — Have you been to China?
\end{quote}

\begin{quote}
\textbf{B:} \zh{去过} \zh{我最喜欢中国菜}
\end{quote}

\begin{quote}
\emph{qù guo, wǒ zuì xǐhuan Zhōngguó cài} — I have. I like Chinese food best.
\end{quote}

\zh{当} answers \zh{当什么}; \zh{因为} answers \zh{为什么}; and B's last line looks back with \zh{过} before it looks forward again.

\begin{tcolorbox}[breakable,skin=enhanced,colback=orange!6,colframe=orange!45!black,boxrule=0.5pt,arc=1mm,left=6pt,right=6pt,top=4pt,bottom=4pt,fonttitle=\bfseries,title={Read this}]
\emph{Read it:} the exchange through once without stopping, then answer the question below
\end{tcolorbox}

\begin{grammarlens}[title={What you've built — looking ahead}]
What you want, what you dream of, what you aim at, and what you like best:

\begin{itemize}
\raggedright
  \item \textbf{\zh{以后我想当医生}} \emph{yǐhòu wǒ xiǎng dāng yīshēng} — a wish for later.
  \item \emph{wǒ de mèngxiǎng shì …} — a dream.
  \item \emph{wǒ de mùbiāo shì …} — a goal.
  \item \textbf{\zh{我最喜欢中国菜}} \emph{wǒ zuì xǐhuan Zhōngguó cài} — the best of all.
\end{itemize}

\textbf{Listen.} \emph{Yǐhòu nǐ xiǎng zuò shénme gōngzuò?}

\textbf{Speak.} Say what you want to be or do later, and one goal on the way.

\emph{Write it:} \zh{当} — the top of \zh{常}, then the foot of \zh{雪}
\end{grammarlens}

\subsection*{Your turn}
\begin{itemize}
\raggedright
  \item \emph{Say it:} \zh{我去过中国} \zh{以后我想在中国工作} — \emph{wǒ qù guo Zhōngguó, yǐhòu wǒ xiǎng zài Zhōngguó gōngzuò} — I've been to China, and later on I want to work there.
  \item \emph{Say it:} \zh{以后你想当什么} — \emph{yǐhòu nǐ xiǎng dāng shénme} — What do you want to be later on?
  \item \emph{Say it:} the whole exchange, both parts
  \item \emph{Write it:} \zh{我想当医生}
\end{itemize}

\begin{tcolorbox}[breakable,colback=teal!4,colframe=teal!35!black,title={Before you move on}]
Say you want to be a doctor. (\textbf{\zh{我想当医生}}.) Ask what someone wants to be later on. (\textbf{\zh{以后你想当什么}}.) Say what you like best. (\textbf{\zh{我最喜欢}…}.)
\end{tcolorbox}
